\documentclass[11pt]{article}
\usepackage[margin=0.9in]{geometry}
\usepackage{amsmath,amssymb,amsthm,booktabs}
\usepackage[T1]{fontenc}
\usepackage{lmodern,microtype}
\usepackage[colorlinks=true,linkcolor=blue,urlcolor=blue]{hyperref}
\newtheorem{theorem}{Theorem}[section]
\newtheorem{lemma}[theorem]{Lemma}
\newtheorem{proposition}[theorem]{Proposition}
\newtheorem{corollary}[theorem]{Corollary}
\theoremstyle{remark}
\newtheorem{remark}[theorem]{Remark}
\newcommand{\R}{\mathbb R}
\newcommand{\C}{\mathbb C}
\newcommand{\norm}[1]{\left\|#1\right\|}
\newcommand{\diag}{\operatorname{diag}}
\newcommand{\ran}{\operatorname{ran}}
\newcommand{\HH}{\mathcal H}
\title{Toward a four-dimensional Yang--Mills construction\\
\large Notebook 11: the full flat constraint-current family}
\author{Research working document}
\date{16 October 2026}
\begin{document}
\maketitle
\begin{abstract}
We bound the nonlinear constraint-current term omitted from the
auxiliary family of Notebook 09. The key estimate follows the actual
corrected averaging operation through fixed-size blocking steps.
Its quadratic contribution at physical spacing $h$ carries a factor
$h$. Flat averages and their physical adjoints are contractions, so
these contributions sum geometrically with a bound independent of
refinement. This yields both a physical Hilbert form bound and a
supremum bound for the constraint vertex. The former gives a uniform
complex current neighborhood for the full homogeneous flat source
inverse, its coarse precision and its extension. The source coordinate
match, first derivative and scope limitations are explicit.
\end{abstract}

\noindent\textbf{Status.}
Research proofs and deposited source identifications checked by separate
internal derivations and audits. No external certification or novelty
claim is made. The result is for $M=L^j$ with a fixed odd $L\ge3$,
which is the refinement sequence used by the blocking construction.
The broader odd-$M$ flat Fourier estimates of Notebook 09 remain
available as inputs, but are not a claim about an arbitrary nonlinear
single-step average of ratio $M$.

\section{Reference operators and the norm}

Write $D$ for the fine scalar gradient, $d_\eta$ for the fine curl,
$Q=Q_M^{\rm st}$ for the normalized straight average, and $Q_0$ for
the scalar block average. Retain the matched source operators
\[
 G'=(D^*D+Q_0^*a'Q_0)^{-1},\quad
 R=I-G'Q_0^*(Q_0G'^2Q_0^*)^{-1}Q_0G',\quad
 \Lambda=G'RD^*,
\]
\[
 \Delta_0=d_\eta^*d_\eta+DRD^*+Q^*aQ,\quad
 G_0=\Delta_0^{-1},\quad q_0=(QG_0Q^*)^{-1},\quad
 E_0=G_0Q^*q_0.
\]
Notebook 09 proves the required inverse, scalar-gauge and auxiliary
current estimates in physical cell Hilbert norms, uniformly in
refinement and compatible periodic volume. The background is fixed
at the identity throughout this notebook. The source value is $a=1$
at unit final block scale; the distinct scalar penalty $a'$ defines
$R$ and is not substituted for $a$.

For fine or coarse cell blocks use
\[
 \norm{T}_\kappa=\max\left\{
  \sup_x\sum_y e^{\kappa d(x,y)}\norm{T(x,y)},
  \sup_y\sum_x e^{\kappa d(x,y)}\norm{T(x,y)}\right\}.
\]
Internal Hilbert norms use physical fine volume $M^{-d}$; coarse
cells use counting measure. Distance is coarse-cell $\ell^1$ distance
or its periodic counterpart. This rectangular operator algebra is
submultiplicative and invariant under the fixed bilinear transpose.
Choose the same small positive $\kappa$ as in Notebook 09.
The color fiber is fixed; its finite constants are included in those
below. Suprema of fine fields mean the supremum of the fixed norm
on each positively oriented bond.
% Notebook 11 body fragment; no preamble.
% Uses amsmath, amssymb, amsthm, standard theorem environments and \norm.
% Every label has the unique prefix fc-.

\section{The separate constraint-current term at flat background}
\label{fc-section}

Fix an odd one-step ratio $L\ge3$ and let $M=L^j$.
The fine spacing is $\eta=M^{-1}$ and the final physical block
scale is one. We use the corrected nonlinear averaging operation
whose flat linearization is $Q_M^{\mathrm{st}}$, the matched scalar
source projection $R$, and the flat inverse and auxiliary-current
theorems already proved. The homogeneous source penalty is $a=1$;
the same conclusions hold for $a$ in a fixed compact positive
interval. All pairings use physical lattice volume. At complex
current, transpose means the fixed bilinear transpose.

The full current inverse in deposited Raw I:13733--13739 is
\begin{equation}
 G_1^J=[\Delta_0-\Theta'_J-\Pi^T\Theta_J\Pi]^{-1},
 \qquad \Pi=I-D\Lambda,
 \label{fc-full-inverse-definition}
\end{equation}
where $\Lambda=G'RD^*$ is fixed at the flat background. The auxiliary
family $G_\pi^J,q_J,E_J$ uses only $\Theta'_J$ and is already
controlled. At final scale one $E_J=H^J$, and the separate term is
\begin{equation}
 \langle A,\Theta_J B\rangle_\eta
   =2\langle C_j^{(2)}(A,B),w_J\rangle_1,
 \qquad w_J=E_J^T J.
 \label{fc-constraint-form}
\end{equation}
Here $C_j^{(2)}$ is the symmetric bilinear quadratic coefficient,
not the full second derivative. The argument below controls this
actual coefficient by composing fixed $L$-step jets. It does not
replace the nonlinear iterate by a newly chosen single $M$-step
rooted average.

\subsection{The corrected coordinate and its one-step jet}
\label{fc-coordinate-section}

The source identifications used here are explicit in the deposit.
Raw I:24320--24325 composes the corrected double-bar averaging
operation itself. I:24344--24369 gives its logarithmic one-step
map and its linear term. I:26985--26988 identifies the tower
$C_j^{(2)}$ with the quadratic term of that nonlinear average.
The scale $\ell$ in the source is the retained physical block
length, hence $\ell=1$ here; it is not $\eta$.

There is also an explicit input-coordinate bridge. Put
\[
 \phi(X;c)=\Phi^{\flat}(1,X;c),\qquad
 \alpha_y(X)=L^{-d}\sum_{x\in B(y)}
                  \Phi^{\mathfrak g}(1,X;y,x).
\]
The maps on the right are defined using $e^{iX}V$ at
I:46965--46973. At $V=1$, I:47076--47087, explicitly incorporating
CMP98 (120), reads
\[
 e^{i\phi(X;c)}=e^{i\alpha_{c_-}(X)}e^{iF_L(X;c)}
                                  e^{-i\alpha_{c_+}(X)}.
\]
Thus the corrected unit-angle map on the identity logarithm branch is
\begin{equation}
 F_L(X;c)=\frac1i\log\!\left[
 e^{-i\alpha_{c_-}(X)}e^{i\phi(X;c)}e^{i\alpha_{c_+}(X)}\right].
 \label{fc-corrected-map}
\end{equation}
Write subscripts $1,2$ for homogeneous Taylor coefficients and
abbreviate $\alpha_{c_\pm}$ by $\alpha_\pm$. BCH gives
\begin{equation}
 \begin{split}
 C_L(X,X)(c)={}&\phi_2-\alpha_{-,2}+\alpha_{+,2}\\
 &+\frac i2\bigl(-[\alpha_{-,1},\phi_1]
                 -[\alpha_{-,1},\alpha_{+,1}]
                 +[\phi_1,\alpha_{+,1}]\bigr).
 \end{split}
 \label{fc-corrected-quadratic}
\end{equation}
Polarization defines $C_L(A,B)$. In particular, the later
unit-lattice vertex named $\Phi^{\flat}_2$ is not substituted for
\eqref{fc-corrected-quadratic}; the endpoint contour and BCH terms
are retained.

At flat background, the explicit linearization I:24357--24369 has
identity transports, $Y=Y_x=0$ and $g(0)=1$. Its correction terms
vanish, leaving
\begin{equation}
 F_L(X)=LQ_L^{\mathrm{st}}X+C_L(X,X)+O(X^3).
 \label{fc-unit-jet}
\end{equation}
Every intermediate reference background is again the identity, so
the same one-step jet is used at every scale.

For fixed $L$, the maps in \eqref{fc-corrected-map} are finite
analytic products, means and logarithms near the identity.
Consequently there is a finite constant $K_L$ such that
\begin{equation}
 |C_L(A,B)(c)|\le K_L
       \max_{b\in\mathcal D_L(c)}|A(b)|
       \max_{b\in\mathcal D_L(c)}|B(b)|,
 \label{fc-one-step-coefficient}
\end{equation}
where $\mathcal D_L(c)$ is the fixed finite set of input bonds read
by the one-step map. The deposited one-step remainder bound
I:24349--24353 also supplies such a constant by taking its quadratic
coefficient. No bound uniform in $L\to\infty$ is needed.

If the fine physical spacing is $h$, then the input angle is $hA$
and the output potential is divided by $Lh$. Therefore
\begin{equation}
 \mathcal F_h(A)=(Lh)^{-1}F_L(hA)
    =Q_hA+\mathcal B_h(A,A)+O(A^3),
 \qquad \mathcal B_h(A,B)=\frac hL C_L(A,B).
 \label{fc-physical-jet}
\end{equation}
The factor $h$ will make the sum over scales convergent.

\subsection{Paired form bounds and the geometric scale sum}
\label{fc-geometric-section}

Let $N_L$ bound the number of output bonds $c$ whose set
$\mathcal D_L(c)$ contains any given input bond. This is a fixed
geometric constant, with covering-space multiplicities on short
tori. Set $k_L=K_LN_LL^{d-1}$.

\begin{lemma}[The one-step physical form]
\label{fc-one-step-form}
For bounded coarse weight $w$,
\begin{equation}
 |\langle w,\mathcal B_h(A,B)\rangle_{Lh}|
 \le k_Lh\norm{w}_\infty\norm{A}_{2,h}\norm{B}_{2,h}.
 \label{fc-one-step-form-bound}
\end{equation}
The flat straight average and its physical adjoint satisfy
\begin{equation}
 \norm{Q_hA}_{2,Lh}\le\norm{A}_{2,h},\qquad
 \norm{Q_h^*w}_\infty\le\norm{w}_\infty.
 \label{fc-contractions}
\end{equation}
\end{lemma}

\begin{proof}
The coarse pairing contributes $(Lh)^d$, while
\eqref{fc-physical-jet} contributes $h/L$. Hence its absolute value
is at most
\[
 K_LL^{d-1}h^{d+1}\norm{w}_\infty
 \sum_c
 \Big(\sum_{b\in\mathcal D_L(c)}|A(b)|^2\Big)^{1/2}
 \Big(\sum_{b\in\mathcal D_L(c)}|B(b)|^2\Big)^{1/2}.
\]
Cauchy--Schwarz over $c$, the incidence bound $N_L$, and the fine
volume factor $h^d$ prove \eqref{fc-one-step-form-bound}.

For a coarse bond in direction $\mu$,
\[
 (Q_h A)_\mu(y)=L^{-d-1}\sum_{x\in B(y)}\sum_{r=0}^{L-1}
                         A_\mu(x+rh e_\mu).
\]
The coefficients are nonnegative and each row sums to one. For
each fine bond there are exactly $L$ string offsets, counting
multiplicity, and each fixes one starting point and block center.
Its column sum is thus $L^{-d}$. The physical adjoint multiplies
this by the coarse/fine volume ratio $L^d$. Therefore both $Q_h$
and $Q_h^*$ have nonnegative rows summing to one. Jensen's
inequality proves the Hilbert contraction, and the adjoint row
identity proves the supremum contraction. This counting also
applies on short periodic tori.
\end{proof}

Set $h_t=L^{t-j}$ for $0\le t\le j$. Let $Q_t$ act from spacing
$h_t$ to $h_{t+1}$ and put
\[
 Q_{r:s}=Q_{s-1}\cdots Q_r\quad(r<s),\qquad Q_{r:r}=I.
\]

\begin{theorem}[Uniform quadratic coefficient of the corrected iterate]
\label{fc-iterated-jet-theorem}
The actual quadratic coefficient has the exact expansion
\begin{equation}
 C_j^{(2)}(A,B)=\sum_{t=0}^{j-1}Q_{t+1:j}
       \mathcal B_{h_t}(Q_{0:t}A,Q_{0:t}B).
 \label{fc-jet-chain}
\end{equation}
With $k_*=k_L/(L-1)$,
\begin{equation}
 |\langle w,C_j^{(2)}(A,B)\rangle_1|
 \le k_*\norm{w}_\infty\norm{A}_{2,\eta}\norm{B}_{2,\eta},
 \label{fc-uniform-jet-form}
\end{equation}
uniformly in $j$, refinement and compatible volume.
\end{theorem}

\begin{proof}
If a partial composition has linear part $P$ and quadratic
coefficient $C$, the next step replaces the latter by
$Q_tC+\mathcal B_{h_t}(P,P)$. This proves
\eqref{fc-jet-chain} inductively, with no factor two because these
are homogeneous coefficients. Pair each summand with $w$, move
the suffix average to $w$, and apply
\eqref{fc-one-step-form-bound} and \eqref{fc-contractions}.
The result is bounded by
\[
 k_L\norm{w}_\infty\norm{A}_{2,\eta}\norm{B}_{2,\eta}
       \sum_{t=0}^{j-1}h_t,
 \qquad
 \sum_{t=0}^{j-1}h_t=\frac{1-L^{-j}}{L-1}.
\]
This proves \eqref{fc-uniform-jet-form}; for $j=0$ the coefficient
is zero. Computing the finite Taylor jet at zero does not require
a common nonlinear neighborhood for the entire iterated map.
\end{proof}

\subsection{The vertex as a local Hilbert and supremum operator}
\label{fc-vertex-norm-section}

Define the symmetric fine operator $T_w$ by
\begin{equation}
 \langle A,T_wB\rangle_\eta
       =2\langle C_j^{(2)}(A,B),w\rangle_1.
 \label{fc-T-definition}
\end{equation}
The previous theorem gives
$\norm{T_w}_{2\to2}\le2k_*\norm{w}_\infty$.
The final nonlinear average reads a fixed number of final unit
cells. The exact source statement is two final blocks
(I:24332--24336); alternatively, a one-step physical radius
$\rho_Lh$ sums to at most $\rho_L/(L-1)$. Thus $T_w(x,y)$ vanishes
when the unit-cell distance exceeds a fixed $R_*$, independent of
$j$. With $v(R_*)=\#\{z\in\mathbb Z^d:|z|_1\le R_*\}$,
\begin{equation}
 \norm{T_w}_\kappa\le
  2k_*v(R_*)e^{\kappa R_*}\norm{w}_\infty
  =:\tau_\kappa\norm{w}_\infty.
 \label{fc-T-weighted}
\end{equation}
Here the weighted norm uses physical cell Hilbert block norms.
Each block norm is bounded by the global Hilbert norm and only a
fixed number of cells is summed. There is no fine-cell dimension
factor.

\begin{lemma}[A stronger bound for the constraint vertex itself]
\label{fc-T-supremum}
The operator $T_w$ is also uniformly bounded on fine
$\ell^\infty$, with fixed coarse-cell range and exponential cell
weights. This assertion does not require a supremum bound for
$\Pi$.
\end{lemma}

\begin{proof}
Write the one-step coefficient in local fine coordinates as
\[
 C_L(A,B)(c)=\sum_{b,b'\in\mathcal D_L(c)}
                    c_{c;b,b'}(A(b),B(b')).
\]
Let $K_L^{\rm coef}$ bound these finitely many bilinear coefficient
norms and $m_L=\max_c|\mathcal D_L(c)|$. The one-step fine
operator $T_w^h$ has absolute row sum at most
\begin{equation}
 2L^{d-1}hN_Lm_LK_L^{\rm coef}\norm{w}_\infty.
 \label{fc-one-step-row}
\end{equation}
The factors are the physical volume ratio $L^d$, the quadratic
factor $h/L$, and the factor two in
\eqref{fc-T-definition}. Each row meets at most $N_L$ outputs and
$m_L$ second input bonds per output. Transpose symmetry gives the
corresponding absolute column bound.

The exact multilevel operator identity is
\begin{equation}
 T_w=\sum_{t=0}^{j-1}Q_{0:t}^*
           T^{h_t}_{Q_{t+1:j}^*w}Q_{0:t}.
 \label{fc-T-operator-chain}
\end{equation}
Both average factors and the backward weight are supremum
contractions. Hence
\begin{equation}
 \norm{T_w}_{\infty\to\infty}
 \le \frac{2L^{d-1}N_Lm_LK_L^{\rm coef}}{L-1}
                         \norm{w}_\infty.
 \label{fc-T-supremum-bound}
\end{equation}
Its fixed final range gives weighted cell-block bounds with
supremum fibers at another fixed range cost. The physical weights
and all periodic multiplicities were already included in
\eqref{fc-one-step-row}.
\end{proof}

\subsection{The actual coefficient weight and the full inverse}
\label{fc-full-neighborhood-section}

Use the auxiliary theorem's constants
\[
 \norm{G_0}_\kappa\le g,\qquad
 \norm{Q}_\kappa,\norm{Q^*}_\kappa\le b,\qquad
 \norm{q_0}_\kappa\le q,
 \qquad c=gb^2q,\quad e=gbq.
\]
On its current ball of radius $\varepsilon_{\rm aux}$,
$\norm{E_J}_\kappa\le4e$. A bounded fine current has physical
Hilbert norm at most $\sqrt d\norm{J}_\infty$ in each unit cell.
The transpose rectangular Schur estimate therefore proves
\begin{equation}
 \norm{w_J}_\infty=\norm{E_J^TJ}_\infty
       \le4\sqrt d\,e\norm{J}_\infty.
 \label{fc-weight-bound}
\end{equation}
This row sum is absolutely convergent also for globally bounded,
non-square-integrable currents in infinite volume.
The map $J\mapsto E_J^TJ$ is analytic into coarse
$\ell^\infty$, because evaluation is a bounded bilinear map in
these norms. Consequently $\Theta_J=T_{w_J}$ is analytic and
\begin{equation}
 \norm{\Theta_J}_\kappa
       \le4\sqrt d\,e\tau_\kappa\norm{J}_\infty.
 \label{fc-Theta-bound}
\end{equation}
It has fixed range in its potential arguments, although its
coefficient weight depends nonlocally on the current.

Let $p=\norm{\Pi}_\kappa<\infty$, using the proved Hilbert bound
on $D\Lambda$, and let $\theta$ be the auxiliary constant
$\norm{\Theta'_J}_\kappa\le\theta\norm{J}_\infty$. Put
\begin{equation}
 \xi=\theta+4p^2\sqrt d\,e\tau_\kappa,
 \qquad
 \varepsilon_1=\min\left\{\varepsilon_{\rm aux},
              \frac1{4g\xi(1+c)}\right\}.
 \label{fc-full-radius}
\end{equation}
The vanishing-perturbation case has the evident interpretation.

\begin{theorem}[The full homogeneous flat current family]
\label{fc-full-family-theorem}
For $\norm{J}_\infty<\varepsilon_1$, the actual full flat operators
\begin{equation}
 G_1^J=[\Delta_0-\Psi_J]^{-1},\quad
 q_1^J=(QG_1^JQ^*)^{-1},\quad E_1^J=G_1^JQ^*q_1^J,
 \qquad \Psi_J=\Theta'_J+\Pi^T\Theta_J\Pi
 \label{fc-full-family}
\end{equation}
are analytic in the weighted physical-cell Hilbert algebra and obey
\begin{equation}
 \norm{G_1^J}_\kappa\le2g,\qquad
 \norm{q_1^J}_\kappa\le2q,\qquad
 \norm{E_1^J}_\kappa\le4e,
 \qquad QE_1^J=I.
 \label{fc-full-family-bounds}
\end{equation}
The constants are uniform in $M=L^j$, compatible volume and the
specified compact positive penalty interval. Fine derivative
sandwiches of $G_1^J$ of total order at most two have the same kind
of uniform weighted Hilbert bounds.
\end{theorem}

\begin{proof}
Equations \eqref{fc-Theta-bound} and \eqref{fc-full-radius} give
$\norm{\Psi_J}_\kappa\le\xi\norm{J}_\infty$ and
$\norm{G_0\Psi_J}_\kappa<1/4$. Thus
\[
 G_1^J=(I-G_0\Psi_J)^{-1}G_0,
 \qquad
 \norm{G_1^J-G_0}_\kappa\le2g^2\xi\norm{J}_\infty.
\]
The coarse inverse exists because
\[
 \norm{q_0 Q(G_1^J-G_0)Q^*}_\kappa
 \le2qb^2g^2\xi\norm{J}_\infty<\frac12.
\]
The second Neumann series gives the bounds in
\eqref{fc-full-family-bounds}. Analyticity follows from the two
norm-convergent inverse series and the already proved analyticity
of $\Psi_J$.

For a prescribed fine-difference sandwich, the resolvent identity is
\[
 \begin{split}
 \mathcal D_\alpha G_1^J\mathcal D_\beta
 ={}&\mathcal D_\alpha G_0\mathcal D_\beta\\
 &+(\mathcal D_\alpha G_0)\Psi_J
       (I-G_0\Psi_J)^{-1}G_0\mathcal D_\beta.
 \end{split}
\]
All needed reference sandwiches and one-sided factors are bounded
for $\alpha+\beta\le2$. Derivatives are placed on the outside
propagators; none acts on $J$.
\end{proof}

\begin{corollary}[Factorial derivatives on smaller current balls]
\label{fc-factorial-corollary}
Let $\mathcal T_J$ denote $G_1^J$, $q_1^J$ or $E_1^J$, with
corresponding uniform bound $M_{\mathcal T}=2g$, $2q$ or $4e$.
For $r=\varepsilon_1-\norm{J}_\infty>0$ and $n\ge1$,
\[
 \norm{D^n\mathcal T_J[K_1,\ldots,K_n]}_\kappa
 \le M_{\mathcal T}\,n!\left(\frac{e_{\rm num}}r\right)^n
                \prod_{i=1}^n\norm{K_i}_\infty,
 \qquad e_{\rm num}=\exp(1).
\]
Here $e_{\rm num}$ is Euler's number, distinct from the extension
constant $e=gbq$.
\end{corollary}
\begin{proof}
Apply the multivariable Cauchy estimate to
$\mathcal T_{J+\sum_i z_i K_i}$ on polydiscs of radii approaching
$r/(n\norm{K_i}_\infty)$. The coefficient containing every $z_i$
once is the displayed mixed derivative. Its norm is bounded by
$M_{\mathcal T}n^nr^{-n}\prod_i\norm{K_i}_\infty$.
Use $n^n\le e_{\rm num}^n n!$. Zero directions follow by multilinearity.
\end{proof}
\subsection{First response, source geometry, and limitations}
\label{fc-response-scope}

The exact first-order constraint form is
\begin{equation}
 \langle A,\Theta_K^{[1]}B\rangle_\eta
    =2\langle C_j^{(2)}(A,B),E_0^T K\rangle_1.
 \label{fc-first-form}
\end{equation}
The derivative of $E_J$ is multiplied by the zero base current and
does not enter. Since $\Pi E_0=E_0$ and the auxiliary first
coarse-precision derivative vanishes,
\begin{equation}
 Dq_1^0[K]=-E_0^T\Theta_K^{[1]}E_0.
 \label{fc-first-precision}
\end{equation}
Equivalently, its bilinear coarse form is
$-2\langle C_j^{(2)}(E_0B,E_0C),E_0^TK\rangle_1$.
This term is not forced to vanish by the auxiliary cancellation.

The full perturbation $\Psi_J$ is nonlinear in $J$. Its derivatives
therefore do not obey the auxiliary ordered-chain formulas without
additional terms. Uniform analytic operator bounds give factorial
derivative estimates on smaller current balls. Spatially, already
\eqref{fc-first-form} may transport a current from cell $x$ to a
coefficient cell $y$ through $E_0^T$, and then branch along two
extension legs to cells $b,c$. Its natural estimate is a sum of
Y-shaped trees over $y$. An auxiliary ordered path
$b\to x\to c$ at the same exponential rate must not be inferred.

There is also a separate pointwise distinction. Lemma
\ref{fc-T-supremum} bounds the constraint vertex itself on fine
$\ell^\infty$. It does not assert that the order-zero projection
$\Pi=I-D\Lambda$ is uniformly bounded in that norm, so it does not
by itself prove a supremum bound for $\Pi^T\Theta_J\Pi$ or for the
full current inverse. The full-family theorem above uses physical
cell Hilbert norms throughout.

A macroscopic path integral need not be bounded uniformly on fine
physical $L^2$: a line of fixed physical length has trace norm of
order $\eta^{(1-d)/2}$. The present proof does not invoke such a
trace. It confines every path coefficient to the fixed $L$-step
geometry, retains the physical factor $h_t$, and uses contractions
before summing the geometric series. The result therefore does not
depend on treating two nonlinear averages with the same linear
part as the same map.

The mathematical estimates have separate internal checks. The
operator identification uses the exact deposited coordinate,
composition and coefficient formulas cited above; the original
CMP98 full text has not been independently reverified in this
calculation. The result concerns the homogeneous flat source
family. It does not establish nonflat current stability, retained
multiscale domains, the nonlinear minimizer beyond the coefficient
used here, or an interacting Yang--Mills continuum limit.

\begin{thebibliography}{9}
\bibitem{repo}
\href{https://github.com/ethanwsussman/Yang--Mills-LLM-experiment/tree/e392ef45013a581664bab76b90867170b30fe10e}
{Pinned Yang--Mills-LLM-experiment account}, commit
\texttt{e392ef45013a581664bab76b90867170b30fe10e}.
The exact deposited passages used for the coordinate and operator
identifications are listed in the text.
\bibitem{balaban98}
T. Ba\l aban, \emph{Averaging operations for lattice gauge theories},
Commun. Math. Phys. 98 (1985), 17--51.
\href{https://doi.org/10.1007/BF01211042}{Publisher record}.
This calculation uses the pinned deposited transcriptions; it does not
claim a new full-text verification of the original paper.
\end{thebibliography}
\end{document}